\documentclass[10pt,a4paper]{amsart}
\usepackage{amsmath,amssymb,amsthm,graphicx}
% Retained preamble for the Stacks Project chapter snapshot in this directory.
% The chapter's own preamble.tex is a copy of the shared Stacks Project preamble; the block
% below is sliced MECHANICALLY, line for line, from it. Nothing was added, removed or
% redefined; the wrapper class, the xr-hyper/hyperref wiring and the externaldocument list
% are omitted because the assembler supplies the document class.
%
% lines taken from preamble.tex: 3, 5, 9, 11, 12, 156, 157, 158, ... (51 lines total)
%
% package coverage restored from the chapter's own preamble (2026-10-13): xy, xy-2cell

\usepackage[all]{xy}
\xyoption{2cell}
\UseAllTwocells

\newenvironment{reference}{\comment}{\endcomment}
\newenvironment{slogan}{\comment}{\endcomment}
\newenvironment{history}{\comment}{\endcomment}
\theoremstyle{plain}
\newtheorem{theorem}[subsection]{Theorem}
\newtheorem{proposition}[subsection]{Proposition}
\newtheorem{lemma}[subsection]{Lemma}
\theoremstyle{definition}
\newtheorem{definition}[subsection]{Definition}
\newtheorem{example}[subsection]{Example}
\newtheorem{exercise}[subsection]{Exercise}
\newtheorem{situation}[subsection]{Situation}
\theoremstyle{remark}
\newtheorem{remark}[subsection]{Remark}
\newtheorem{remarks}[subsection]{Remarks}
\numberwithin{equation}{subsection}
\def\lim{\mathop{\mathrm{lim}}\nolimits}
\def\colim{\mathop{\mathrm{colim}}\nolimits}
\def\Spec{\mathop{\mathrm{Spec}}}
\def\Hom{\mathop{\mathrm{Hom}}\nolimits}
\def\Ext{\mathop{\mathrm{Ext}}\nolimits}
\def\SheafHom{\mathop{\mathcal{H}\!\mathit{om}}\nolimits}
\def\SheafExt{\mathop{\mathcal{E}\!\mathit{xt}}\nolimits}
\def\Sch{\mathit{Sch}}
\def\Mor{\mathop{\mathrm{Mor}}\nolimits}
\def\Ob{\mathop{\mathrm{Ob}}\nolimits}
\def\Sh{\mathop{\mathit{Sh}}\nolimits}
\def\NL{\mathop{N\!L}\nolimits}
\def\CH{\mathop{\mathrm{CH}}\nolimits}
\def\proetale{{pro\text{-}\acute{e}tale}}
\def\etale{{\acute{e}tale}}
\def\QCoh{\mathit{QCoh}}
\def\Ker{\mathop{\mathrm{Ker}}}
\def\Im{\mathop{\mathrm{Im}}}
\def\Coker{\mathop{\mathrm{Coker}}}
\def\Coim{\mathop{\mathrm{Coim}}}
\DeclareMathSymbol{\boxtimes}{\mathbin}{AMSa}{"02}
\def\QCohstack{\mathcal{QC}\!\mathit{oh}}
\def\Cohstack{\mathcal{C}\!\mathit{oh}}
\def\Spacesstack{\mathcal{S}\!\mathit{paces}}
\def\Quotfunctor{\mathrm{Quot}}
\def\Hilbfunctor{\mathrm{Hilb}}
\def\Curvesstack{\mathcal{C}\!\mathit{urves}}
\def\Polarizedstack{\mathcal{P}\!\mathit{olarized}}
\def\Complexesstack{\mathcal{C}\!\mathit{omplexes}}
\def\Pic{\mathop{\mathrm{Pic}}\nolimits}
\def\Picardstack{\mathcal{P}\!\mathit{ic}}
\def\Picardfunctor{\mathrm{Pic}}
\def\Deformationcategory{\mathcal{D}\!\mathit{ef}}
\usepackage{cleveref}
\crefname{originallemma}{original lemma}{original lemmas}
\Crefname{originallemma}{Original lemma}{Original lemmas}
\makeatletter
\newlabel{dga-lemma-id-cone-null}{{09QK}{source}{}{sourcepointer-1}{}}
\newlabel{dga-lemma-id-cone-null@cref}{{[originallemma][0][]09QK}{[1][0][]source}{}{}{}}
\makeatother
\crefname{originallemma}{original lemma}{original lemmas}
\Crefname{originallemma}{Original lemma}{Original lemmas}
\makeatletter
\newlabel{homology-lemma-ML-over-ordinals}{{0AAT}{source}{}{sourcepointer-2}{}}
\newlabel{homology-lemma-ML-over-ordinals@cref}{{[originallemma][0][]0AAT}{[1][0][]source}{}{}{}}
\makeatother
\crefname{originallemma}{original lemma}{original lemmas}
\Crefname{originallemma}{Original lemma}{Original lemmas}
\makeatletter
\newlabel{injectives-lemma-characterize-K-injective}{{079L}{source}{}{sourcepointer-3}{}}
\newlabel{injectives-lemma-characterize-K-injective@cref}{{[originallemma][0][]079L}{[1][0][]source}{}{}{}}
\makeatother
\crefname{originallemma}{original lemma}{original lemmas}
\Crefname{originallemma}{Original lemma}{Original lemmas}
\makeatletter
\newlabel{lemma-characterize-graded-injectives-in-dg}{{0FSS}{source}{}{sourcepointer-4}{}}
\newlabel{lemma-characterize-graded-injectives-in-dg@cref}{{[originallemma][0][]0FSS}{[1][0][]source}{}{}{}}
\makeatother
\crefname{originallemma}{original lemma}{original lemmas}
\Crefname{originallemma}{Original lemma}{Original lemmas}
\makeatletter
\newlabel{lemma-first-property-dg-injective}{{0FSW}{source}{}{sourcepointer-5}{}}
\newlabel{lemma-first-property-dg-injective@cref}{{[originallemma][0][]0FSW}{[1][0][]source}{}{}{}}
\makeatother
\crefname{originallemma}{original lemma}{original lemmas}
\Crefname{originallemma}{Original lemma}{Original lemmas}
\makeatletter
\newlabel{lemma-small-acyclics}{{0FST}{source}{}{sourcepointer-6}{}}
\newlabel{lemma-small-acyclics@cref}{{[originallemma][0][]0FST}{[1][0][]source}{}{}{}}
\makeatother
\crefname{originallemma}{original lemma}{original lemmas}
\Crefname{originallemma}{Original lemma}{Original lemmas}
\makeatletter
\newlabel{remark-why-graded-injective}{{0FSQ}{source}{}{sourcepointer-7}{}}
\newlabel{remark-why-graded-injective@cref}{{[originallemma][0][]0FSQ}{[1][0][]source}{}{}{}}
\makeatother
\crefname{originallemma}{original lemma}{original lemmas}
\Crefname{originallemma}{Original lemma}{Original lemmas}
\makeatletter
\newlabel{section-conclude-triangulated}{{0FS3}{source}{}{sourcepointer-8}{}}
\newlabel{section-conclude-triangulated@cref}{{[originallemma][0][]0FS3}{[1][0][]source}{}{}{}}
\makeatother
\title{Stacks, Differential Graded Sheaves: lemma-better-set-of-monos}
\author{The Stacks Project Authors}
\date{stacks-project:sdga.tex:local-snapshot-20261012; GNU FDL 1.2 or later}
\begin{document}
\maketitle
\noindent\textit{Editorial note.} Statement span is lines 4061-4083 of sources/stacks-sdga/pinned-original-source.tex; proof span is lines 4085-4181, the complete author proof (95 body lines).\par\medskip
\noindent\textit{Editorial note.} Navigation stubs: 8 cited-but-not-included labels are supplied as non-load-bearing pointers with their pinned tags where the snapshot tag index has them.\par\medskip
\noindent\textit{Editorial note.} extra\_packages is empty on purpose: the retained preamble already loads xy with options and a second declaration clashes in the assembler wrapper.\par\medskip
\noindent\textit{Editorial note.} Text is retained verbatim, including the author's own typos and placeholders.\par\medskip
\noindent\textit{Editorial note.} Queue index 1 of queue\_sdga.json (pb=4158, stmt 4061-4083, proof 4085-4181).\par\medskip
\noindent\textit{Allowed standard prerequisites (editorial).}\begin{itemize}
\item Explicit prerequisite boundary. The cited inputs are prerequisites; the argument in the retained span is not replaced by any cited result.
\end{itemize}\medskip
\ifdefined\hypertarget\hypertarget{sourcepointer-1}{}\fi
\noindent\textit{Source reference.} Original lemma 09QK: sources/stacks-sdga/pinned-original-source.tex, source lines 4181--4181. Cited by the retained statement or proof but not part of the collected spans; supplied only as a navigation pointer. The argument does not depend on its content.\par\medskip
\ifdefined\hypertarget\hypertarget{sourcepointer-2}{}\fi
\noindent\textit{Source reference.} Original lemma 0AAT: sources/stacks-sdga/pinned-original-source.tex, source lines 4181--4181. Cited by the retained statement or proof but not part of the collected spans; supplied only as a navigation pointer. The argument does not depend on its content.\par\medskip
\ifdefined\hypertarget\hypertarget{sourcepointer-3}{}\fi
\noindent\textit{Source reference.} Original lemma 079L: sources/stacks-sdga/pinned-original-source.tex, source lines 4181--4181. Cited by the retained statement or proof but not part of the collected spans; supplied only as a navigation pointer. The argument does not depend on its content.\par\medskip
\ifdefined\hypertarget\hypertarget{sourcepointer-4}{}\fi
\noindent\textit{Source reference.} Original lemma 0FSS: sources/stacks-sdga/pinned-original-source.tex, source lines 3795--3815. Cited by the retained statement or proof but not part of the collected spans; supplied only as a navigation pointer. The argument does not depend on its content.\par\medskip
\ifdefined\hypertarget\hypertarget{sourcepointer-5}{}\fi
\noindent\textit{Source reference.} Original lemma 0FSW: sources/stacks-sdga/pinned-original-source.tex, source lines 3995--4010. Cited by the retained statement or proof but not part of the collected spans; supplied only as a navigation pointer. The argument does not depend on its content.\par\medskip
\ifdefined\hypertarget\hypertarget{sourcepointer-6}{}\fi
\noindent\textit{Source reference.} Original lemma 0FST: sources/stacks-sdga/pinned-original-source.tex, source lines 3852--3860. Cited by the retained statement or proof but not part of the collected spans; supplied only as a navigation pointer. The argument does not depend on its content.\par\medskip
\ifdefined\hypertarget\hypertarget{sourcepointer-7}{}\fi
\noindent\textit{Source reference.} Original lemma 0FSQ: sources/stacks-sdga/pinned-original-source.tex, source lines 3731--3773. Cited by the retained statement or proof but not part of the collected spans; supplied only as a navigation pointer. The argument does not depend on its content.\par\medskip
\ifdefined\hypertarget\hypertarget{sourcepointer-8}{}\fi
\noindent\textit{Source reference.} Original lemma 0FS3: sources/stacks-sdga/pinned-original-source.tex, source lines 2561--2579. Cited by the retained statement or proof but not part of the collected spans; supplied only as a navigation pointer. The argument does not depend on its content.\par\medskip
\begin{lemma}
\label{lemma-better-set-of-monos}
Let $(\mathcal{C}, \mathcal{O})$ be a ringed site.
Let $(\mathcal{A}, \text{d})$
be a sheaf of differential graded algebras on $(\mathcal{C}, \mathcal{O})$.
There exists a set $R$ and for each $r \in R$ an injective map
$\mathcal{M}_r \to \mathcal{M}'_r$ of
acyclic differential graded $\mathcal{A}$-modules
such that for an object $\mathcal{I}$ of $\textit{Mod}(\mathcal{A}, \text{d})$
the following are equivalent
\begin{enumerate}
\item $\mathcal{I}$ is K-injective and graded injective, and
\item for every solid diagram
$$
\xymatrix{
\mathcal{M}_r \ar[r] \ar[d] & \mathcal{I} \\
\mathcal{M}'_r \ar@{..>}[ru]
}
$$
a dotted arrow exists in $\textit{Mod}(\mathcal{A}, \text{d})$
making the diagram commute.
\end{enumerate}
\end{lemma}

\begin{proof}
Let $T$ and $\mathcal{M}_t \to \mathcal{M}'_t$ be as in
Lemma \ref{lemma-characterize-graded-injectives-in-dg}.
Let $S$ and $\mathcal{M}_s$ be as in
Lemma \ref{lemma-small-acyclics}.
Choose an injective map $\mathcal{M}_s \to \mathcal{M}'_s$
of acyclic differential graded $\mathcal{A}$-modules
which is homotopic to zero. This is possible because we
may take $\mathcal{M}'_s$ to be the cone on the identity;
in that case it is even true that the identity on
$\mathcal{M}'_s$ is homotopic to zero, see
Differential Graded Algebra, Lemma \ref{dga-lemma-id-cone-null} which
applies by the discussion in Section \ref{section-conclude-triangulated}.
We claim that $R = T \coprod S$ with the given maps works.

\medskip\noindent
The implication (1) $\Rightarrow$ (2) holds by
Lemma \ref{lemma-first-property-dg-injective}.

\medskip\noindent
Assume (2). First, by Lemma \ref{lemma-characterize-graded-injectives-in-dg}
we see that $\mathcal{I}$ is graded injective. Next, let
$\mathcal{M}$ be an acyclic differential graded $\mathcal{A}$-module.
We have to show that
$$
\Hom_{K(\textit{Mod}(\mathcal{A}, \text{d}))}(\mathcal{M}, \mathcal{I}) = 0
$$
The proof will be exactly the same as the proof of
Injectives, Lemma \ref{injectives-lemma-characterize-K-injective}.

\medskip\noindent
We are going to construct by induction on the ordinal $\alpha$
an acyclic differential graded submodule
$\mathcal{K}_\alpha \subset \mathcal{M}$ as follows.
For $\alpha = 0$ we set $\mathcal{K}_0 = 0$. For $\alpha > 0$
we proceed as follows:
\begin{enumerate}
\item If $\alpha = \beta + 1$ and $\mathcal{K}_\beta = \mathcal{M}$
then we choose $\mathcal{K}_\alpha = \mathcal{K}_\beta$.
\item If $\alpha = \beta + 1$ and $\mathcal{K}_\beta \not = \mathcal{M}$
then $\mathcal{M}/\mathcal{K}_\beta$ is a nonzero acyclic
differential graded $\mathcal{A}$-module.
We choose a differential graded $\mathcal{A}$ submodule
$\mathcal{N}_\alpha \subset \mathcal{M}/\mathcal{K}_\beta$
isomorphic to $\mathcal{M}_s$ for some $s \in S$, see
Lemma \ref{lemma-small-acyclics}.
Finally, we let $\mathcal{K}_\alpha \subset \mathcal{M}$
be the inverse image of $\mathcal{N}_\alpha$.
\item If $\alpha$ is a limit ordinal we set
$\mathcal{K}_\beta = \colim \mathcal{K}_\alpha$.
\end{enumerate}
It is clear that $\mathcal{M} = \mathcal{K}_\alpha$ for a suitably large
ordinal $\alpha$. We will prove that
$$
\Hom_{K(\textit{Mod}(\mathcal{A}, \text{d}))}(\mathcal{K}_\alpha, \mathcal{I})
$$
is zero by transfinite induction on $\alpha$. It holds for $\alpha = 0$
since $\mathcal{K}_0$ is zero. Suppose it holds for $\beta$ and
$\alpha = \beta + 1$. In case (1) of the list above the result is clear.
In case (2) there is a short exact sequence
$$
0 \to \mathcal{K}_\beta \to \mathcal{K}_\alpha \to \mathcal{N}_\alpha \to 0
$$
By Remark \ref{remark-why-graded-injective}
and since we've seen that $\mathcal{I}$ is graded
injective, we obtain an exact sequence
$$
\Hom_{K(\textit{Mod}(\mathcal{A}, \text{d}))}(\mathcal{K}_\beta, \mathcal{I})
\to
\Hom_{K(\textit{Mod}(\mathcal{A}, \text{d}))}(\mathcal{K}_\alpha, \mathcal{I})
\to
\Hom_{K(\textit{Mod}(\mathcal{A}, \text{d}))}(\mathcal{N}_\alpha, \mathcal{I})
$$
By induction the term on the left is zero. By assumption (2)
the term on the right is zero: any map $\mathcal{M}_s \to \mathcal{I}$
factors through $\mathcal{M}'_s$ and hence is homotopic to zero.
Thus the middle group is zero too. Finally, suppose that $\alpha$ is a
limit ordinal. Because we also have
$\mathcal{K}_\alpha = \colim \mathcal{K}_\alpha$ as graded
$\mathcal{A}$-modules we see that
$$
\Hom_{\textit{Mod}^{dg}(\mathcal{A}, \text{d})}
(\mathcal{K}_\alpha, \mathcal{I})
= \lim_{\beta < \alpha}
\Hom_{\textit{Mod}^{dg}(\mathcal{A}, \text{d})}
(\mathcal{K}_\beta, \mathcal{I})
$$
as complexes of abelian groups. The cohomology groups of these
complexes compute morphisms in $K(\textit{Mod}(\mathcal{A}, \text{d}))$
between shifts. The transition maps in the system of complexes
are surjective by Remark \ref{remark-why-graded-injective}
because $\mathcal{I}$ is graded injective.
Moreover, for a limit ordinal $\beta \leq \alpha$
we have equality of limit and value. Thus we may apply
Homology, Lemma \ref{homology-lemma-ML-over-ordinals}
to conclude.
\end{proof}

\end{document}
